\chapter{项目目标、范围与阅读方式}
\section{要解决的问题}
项目的目标是让机器人根据图像和自然语言指令执行操作，而不只是输出文字描述或重复预录轨迹。为此，需要分别处理目标理解、操作位置选择、动作生成、物理执行、完成判断和运行维护。

本报告中，VLM（视觉语言模型）主要指Qwen这类接收图像和文字、输出文字或位置描述的模型；VLA（视觉语言动作模型）主要指\pifive{}、SmolVLA这类直接生成机器人动作的模型。Qwen输出位置后交给规划器执行，是分模块系统；它与直接由动作模型产生关节或末端动作的路径分别介绍。
\section{系统逐步形成的三条路径}
\begin{enumerate}
\item 观察路径：相机图像进入目标检测、深度显示和场景描述模块，供操作人员观察。
\item 策略路径：图像、语言和机器人状态进入动作模型，模型输出动作块，执行后再次观察。
\item 显式定位路径：视觉模型给出二维操作点，深度和标定信息将其转换到机器人坐标，再通过规划器与控制器执行。
\end{enumerate}
三条路径有共享模块，但输入、输出和验收方式不同。早期演示、公开权重复现、自训练模块以及真实硬件试验，在正文中分别标注。
\section{技术篇章与时间线}
本报告按问题组织六个技术篇章，每篇四个小阶段。章节编号体现技术关系，不强行等同于日期先后；例如端侧环境准备与可供性执行开发曾交错进行。早期阶段主要在2026年10月，统一VLA验证、同场景操作与新端侧评测集中在11月。
\begin{table}[htbp]\centering\begin{tabular}{p{.10\linewidth}p{.31\linewidth}p{.49\linewidth}}
\toprule 篇章 & 主题 & 回答的核心问题\\\midrule
一 & 原型与初始基线 & 数据、观察与动作接口是否可用？\\
二 & 从演示到自主操作 & 现有策略能否真正产生任务动作？\\
三 & 分层与视觉检查 & 高层模型在何时提供了额外作用？\\
四 & 目标与移动操作 & 目标能否传到执行端，导航操作能否同场景？\\
五 & 可供性与物理接口 & 预测操作点能否改变实际到位位置？\\
六 & 小模型与端侧部署 & 精度、内存和时延能否同时满足要求？\\\bottomrule
\end{tabular}\caption{技术篇章与关注问题}\end{table}
\section{结果使用规则}
成功率随任务、初态、执行预算和判定条件变化。预接触到位、单物体搬运、放入容器、原LIBERO整任务成功分别统计；开发回合、验证集和测试集分别记录。图像示例用于解释场景与错误，不能替代总体评测。视频随报告保存到英文命名目录，正文给出相对路径。

\section{不同阶段的输入、输出与控制条件}
\begin{longtable}{p{.19\linewidth}p{.23\linewidth}p{.23\linewidth}p{.23\linewidth}}
\toprule 路径 & 输入 & 输出 & 执行和评分\\\midrule\endhead
SO101示教与回放 & 主臂关节、相机 & 从臂关节目标 & 真实舵机；人工观察任务\\
桌面几何与ACT & 彩色／深度图或状态 & 末端动作序列 & 理想位置驱动；仿真判定\\
公开LIBERO策略 & 两路图像、指令、状态 & 连续动作块 & 动力学仿真；官方成功判定\\
Qwen高层 & 场景图、任务或阶段记忆 & 指令、判断或操作点 & 交给既有策略或规划器\\
可供性到位 & 二维点、深度与相机参数 & 世界坐标和关节轨迹 & MoveIt规划；关节力矩控制\\
Jetson视觉服务 & 与PC一致的图像请求 & 点、框等结构化结果 & PC等待返回；单独测时延\\
\bottomrule\caption{系统路径的能力边界}\end{longtable}
\section{真值信息在不同实验中的用途}
仿真可以提供精确状态，但使用方式会影响结论。公开基准的官方成功检查用于评分与原生回合停止；分层策略的中间谓词仅作事后记录，最后一个子目标仍按原生成功或步数上限终止。后续同场景里程碑版本才验证由视觉模型主动停止，二者不能混称。

可供性训练使用分割和位置真值自动标注；三维反投影使用仿真深度，碰撞场景也由仿真构造。因此，操作点由视觉模型预测，并不意味着深度、建图和避障也由模型估计。导航使用已知地图与理想里程计，不属于未知环境中的自主建图定位。报告在对应实验中保留这些设置。
\section{数据划分和计量单位}
一帧图像可能对应把手、碗、放置点和柜子框等多条查询；350帧与1014条查询是同一测试集的两种计数，不能当成两个独立样本集。仿真回合、视觉检查次数、策略调用次数和控制步数也分别计量。延迟中的秒指实际经过的时间；仿真步数除以仿真频率得到模拟时间，模型推理时模拟器暂停的实验不满足真实时间运行的含义。

可供性数据按初态隔离：训练17--49、验证13--16、测试0--12；不同初态仍共享生成动作模板和相机配置，因此只检验场景内变化，不能当作新动作模式或新场景泛化。早期100张检测训练图片中的20张拟合检查仍属于训练数据，正文不会将其称为独立验证。
\section{模型与环境的分开记录}
模型名称相同但量化范围、库版本、提示词和动作接口不同，输出可能变化。10月借用Jetson与11月恢复后的板子属于两套软件栈；前者的量化范围包含视觉线性层，后者保留视觉分支和输出层的较高精度。不同阶段的数字按各自条件报告，不作直接硬件优劣排序。

本报告的固定源版本和运行命令均以项目内记录为准，附录提供阅读入口。代码目录按章节索引，但共享实现保留在原位置，避免为报告复制后产生多套行为不同的实现。报告没有执行新的训练、闭环或硬件测试。
